% Generated draft: requires mathematical and source review.
\documentclass[11pt]{article}
\usepackage[margin=25mm]{geometry}
\usepackage{fontspec}
\setmainfont{Latin Modern Roman}
\usepackage{amsmath,amssymb,mathtools,unicode-math}
\setmathfont{Latin Modern Math}
\usepackage{enumitem,needspace,xurl,hyperref}
\setlength{\parindent}{0pt}
\setlength{\parskip}{5pt}
\setlength{\emergencystretch}{3em}
\newcommand{\N}{\mathbb{N}}
\newcommand{\Z}{\mathbb{Z}}
\newcommand{\Q}{\mathbb{Q}}
\newcommand{\R}{\mathbb{R}}
\newcommand{\C}{\mathbb{C}}
\newcommand{\F}{\mathbb{F}}
\newcommand{\A}{\mathbb{A}}
\newcommand{\PP}{\mathbb{P}}
\newcommand{\E}{\mathbb{E}}
\newcommand{\eps}{\varepsilon}
\newcommand{\dd}{\,\mathrm{d}}
\begin{document}
\section*{UnsolvedMath ID: 11000272}
\textbf{Source problem code:} catalogue entry
\subsection{Definitions and mathematical statement}
Let \(F_n\) be the free group of rank \(n\), and \(\operatorname{Aut}(F_n)\) its automorphism group. A group has Kazhdan's property \((T)\) if every unitary representation with almost invariant unit vectors has a nonzero invariant vector. Question 2.4 asks, for \(n>3\), whether \(\operatorname{Aut}(F_n)\) has property \((T)\). The strict range is \(n\ge4\); the separately discussed rank-three case is not the same question.

The supplied research summary labels the problem solved negatively for every \(n\ge3\), without a primary citation, invoking unspecified unbounded actions, abelianization/free quotients and the surjection onto \(\mathrm{GL}(n,\Z)\). Those phrases do not establish the claim: arbitrary unbounded actions are not the unitary obstruction required by \((T)\), and a quotient map alone is not an obstruction unless the quotient is shown to fail \((T)\). No asserted rank-at-least-four free/abelian quotient is constructed. Thus the detailed negative assessment is preserved as unverified and is not promoted to a theorem.

The source gives the following context. A positive \((T)\) answer would, through Lubotzky--Pak [64], conceptually explain effectiveness of the Product Replacement Algorithm of Leedham-Green and others; no quantitative mixing statement is supplied. It also says the mapping-class analogue was open in the source era.

The source calls \(F_A\) the property that every action on an \(\R\)-tree has a fixed point and reports it for \(\mathrm{GL}(n,\Z)\), \(\operatorname{Aut}(F_n)\), \(\operatorname{Out}(F_n)\) when \(n\ge3\), and genus-at-least-three mapping class groups, reference [28]. Standard Serre \(F_A\) concerns simplicial trees, whereas the all-\(\R\)-tree version is commonly denoted \(F_\R\); the intended convention and hypotheses must be checked. The source also states that \((T)\) implies its fixed-point property.

Finally it cites McCool [67] for a finite-index subgroup of \(\operatorname{Aut}(F_3)\) with positive first Betti number and hence a quotient \(\Z\). Translation on a line then gives a fixed-point-free action of that subgroup; since \((T)\) passes to finite-index subgroups, it concludes failure of \((T)\) for \(\operatorname{Aut}(F_3)\). This rank-three argument does not justify the summary's claimed negative conclusion for all larger ranks.
\subsection{Short English statement}
For ranks strictly greater than three, ask whether the free-group automorphism group has property \((T)\). The source separately gives a rank-three obstruction. Its summary claims a negative solution in all ranks without a citation or valid supplied argument, so that status claim remains explicitly unresolved here.
\subsection{Sources}
\begin{itemize}
\item Local UnsolvedMath registry, permanent ID \(11000272\), exact source code \texttt{catalogue entry}, Question 2.4, PDF page 330. No individual proposer or proposed year is supplied.
\item Farb, edited volume \emph{Problems on Mapping Class Groups and Related Topics}: \url{https://www.math.uchicago.edu/~farb/papers/mcgbook.pdf}. The volume has 329 explicitly typeset numbered items across chapters, repeated local numbering distinguished by PDF page.
\item Internal references: Lubotzky--Pak [64] for the Product Replacement motivation; [28] for tree fixed-point properties; McCool [67] for the rank-three finite-index first-Betti result. Full bibliography entries are absent. Leedham-Green and others are credited for the algorithm.
\item Extraction pdf-font-structure, accessed 2026-07-29. Source evidence NEEDS\_REVIEW presents an open problem and requires a current primary-literature audit. Registry status solved and classification SOLVED-IN-LITERATURE follow a summary claiming a negative answer for all \(n\ge3\), but no primary resolution reference is supplied and its reasoning does not establish that rank range. This conflict is retained without guessing a modern answer.
\item Dataset attribution: Ulam AI / UnsolvedMath Contributors, CC BY 4.0. Source rights NEEDS\_REVIEW: publicly accessible author/publisher copy, redistribution terms requiring review.
\end{itemize}
\end{document}
